%! MDraft's built-in LaTeX reference.
%!
%! Opened by Help › LaTeX Reference. Every example outside a verbatim block is
%! compiled by the same LuaLaTeX the preview uses, so what this file says about
%! the syntax is what the application actually renders. A built-in document has
%! no folder of its own, so `\includegraphics` and `\input` are shown as source
%! rather than called with a path that could not resolve.

\documentclass[11pt]{article}
\usepackage{amsmath}
\usepackage{amssymb}
\usepackage{hyperref}
\usepackage{xcolor}
\usepackage[normalem]{ulem}

\title{LaTeX Reference}
\author{MDraft}
\date{}

\begin{document}
\maketitle

\section{Introduction}

This document is a live sample. LaTeX is a typesetting system: the text below is
markup, the lines beginning with a backslash are commands, and LuaLaTeX evaluates
both before it lays out a page. Edit any line and the preview recompiles.

The preview draws the PDF the compiler produced, with the real page size, margins
and layout. Select text in it to copy it, and click it to jump to the source line
that produced it.

\tableofcontents

\section{Headings and text}

A heading is a command: \verb|\section{...}| is level one,
\verb|\subsection{...}| is level two, and \verb|\subsubsection{...}| is level
three. The starred form (\verb|\section*{...}|) is unnumbered. The Outline panel
reads these commands, so the sidebar and the document cannot disagree.

Paragraphs are separated by a blank line. A single newline
continues the same paragraph. Two backslashes force a line break,\\
like this. A blank line starts a new paragraph.

\textbf{Strong}, \emph{emphasis}, \texttt{inline code}, \underline{underline},
\sout{strikethrough}, H\textsubscript{2}O and E = mc\textsuperscript{2} are all
part of the markup. A backslash escapes a literal character: \# \$ \% \& \_ are
text, not syntax.

\section{Mathematics}

Inline mathematics sits between dollar signs: $e^{i\pi} + 1 = 0$. Display
mathematics uses \verb|\[ ... \]|:

\[
  \int_{-\infty}^{\infty} e^{-x^2}\,dx = \sqrt{\pi}
\]

A numbered equation uses the \verb|equation| environment, and a multi-line
alignment uses \verb|align|:

\begin{align}
  (a + b)^2 &= a^2 + 2ab + b^2 \\
  (a - b)^2 &= a^2 - 2ab + b^2
\end{align}

Fractions, radicals, sums and integrals are all ordinary commands:
$\frac{a}{b}$, $\sqrt{x}$, $\sum_{n=1}^{\infty} \frac{1}{n^2}$,
$\lim_{x \to 0} \frac{\sin x}{x}$. The application's own math settings do not
apply here: LaTeX typesets its own mathematics.

\section{Lists and tables}

An unordered list is the \verb|itemize| environment; an ordered list is
\verb|enumerate|. Each entry begins with \verb|\item|.

\begin{itemize}
  \item A first point.
  \item A second point, with \emph{emphasis}.
  \begin{itemize}
    \item A nested point.
  \end{itemize}
\end{itemize}

\begin{enumerate}
  \item A numbered step.
  \item Another step.
\end{enumerate}

A table is a \verb|tabular| with a column specification: \verb|l| and \verb|c|
are left- and centre-aligned columns, and \verb|\hline| draws a rule.

\begin{table}[htbp]
  \centering
  \begin{tabular}{ll}
    \hline
    Command & Purpose \\
    \hline
    \verb|\textbf| & Strong text \\
    \verb|\emph| & Emphasis \\
    \verb|\texttt| & Inline code \\
    \hline
  \end{tabular}
  \caption{A small table.}
\end{table}

\section{Quotes and verbatim}

The \verb|quote| environment indents a passage:

\begin{quote}
  A quotation is set apart from the surrounding text, without quotation marks of
  its own.
\end{quote}

Verbatim text keeps every character exactly as written, which is how code is
shown. The \verb|verbatim| environment is a block; \verb|\verb| is inline.

\begin{verbatim}
for i in range(3):
    print(i)   # % is not a comment here
\end{verbatim}

\section{Links, references and citations}

An external link is \verb|\href{url}{text}|: \href{https://www.latex-project.org/}{the LaTeX project}.
A bare URL is \verb|\url{...}|: \url{https://www.latex-project.org/}.

A label marks a place and a reference points at it. This section is labelled
\verb|sec:links|, and \verb|\ref{sec:links}| prints its number: see
Section~\ref{sec:links}. The \verb|\pageref| command prints the page instead.

A citation reads a bibliography. The engine runs real BibTeX, so a
\verb|\cite{...}| key must exist in the \verb|.bib| file named by
\verb|\bibliography{...}|. Biber and the full \verb|biblatex| feature set are
not available; \verb|biblatex-lite| covers the common cases.

\section{Including other files}

A large document is split with \verb|\input{chapters/intro}| or
\verb|\include{chapters/intro}|.\footnote{The \texttt{.tex} extension is
optional.} Paths are relative to the main file's folder, which is also the
project root: a reference that escapes it is an error rather than a silent
omission.

An image is \verb|\includegraphics{figures/plot}|, optionally with a size:
\verb|\includegraphics[width=0.5\linewidth]{figures/plot}|. The application
inserts this command when you paste or drop an image into a \verb|.tex|
document.

\section{Exporting}

\verb|File > Export > PDF| writes the PDF this preview is showing, byte for byte.
It never starts a compile of its own: if the preview is stale, the row is
disabled and says so. The preview's own \verb|Save PDF| button writes the same
bytes. Native printing is disabled for LaTeX because the PDF already exists.

\section{When something fails}

A compile error appears above the preview with the compiler's own message and
the line it refers to. The previous successful PDF stays on screen, so a
half-finished edit does not blank the page. Fixing the error and waiting a moment
recompiles automatically.

The engine is opt-in. If the preview offers \verb|Enable LuaLaTeX|, the engine
files have not been downloaded yet; about 10 MiB of verified files are fetched
once, and the TeX packages a document needs are fetched on demand and cached
locally. Nothing about your document leaves this device.

\end{document}
